\chapter{Symmetry, Involution, and the Half-Axis}
\label{chap:symmetry}

\section{The conjugate-affine involution}

The functional equation is a holomorphic reflection $s\mapsto1-s$. To express reflection across a vertical line, combine it with complex conjugation.

\begin{definition}[Critical-strip involution]
Define
\begin{equation}
 \J:\C\to\C,\qquad \J(s)=1-\overline{s}.
 \label{eq:J-def}
\end{equation}
\end{definition}
The map $\J$ is an anti-holomorphic, conjugate-affine involution and satisfies $\J^2=\operatorname{id}$. It is not anti-linear as a map of the affine coordinate $s$ because of the constant term $1$. After centering at $1/2$, however, its action becomes anti-linear; this distinction is used explicitly below.

\begin{lemma}[SOH-L003: fixed set]
\label{lem:J-fixed}
The fixed set of $\J$ is exactly the critical line:
\[
 \operatorname{Fix}(\J)=\Crit.
\]
\end{lemma}
\begin{proof}
Write $s=\sigma+it$. Then
\[
 \J(s)=1-\sigma+it.
\]
The equation $\J(s)=s$ is equivalent to $1-\sigma=\sigma$, hence $\sigma=1/2$, while $t$ remains arbitrary.
\end{proof}

The completed zeta function is covariant under this involution:
\begin{equation}
 \xi(\J s)=\overline{\xi(s)}.
 \label{eq:xi-J-covariance}
\end{equation}
Indeed,
\[
 \xi(1-\overline{s})=\xi(\overline{s})=\overline{\xi(s)}.
\]
This equation will later be the model for covariance of a candidate state map.

\section{Why symmetry does not imply the Riemann Hypothesis}

A common logical mistake is to infer that because the zero set is invariant under $\J$, every zero must satisfy $\J(\rho)=\rho$. Set invariance is weaker than pointwise fixation.

\begin{proposition}[Orbit alternative]
Let $Z\subset\Sstrip$ be invariant under conjugation and $s\mapsto1-s$. For every $\rho\in Z$, either:
\begin{enumerate}
  \item $\Re\rho=1/2$, in which case $\J(\rho)=\rho$; or
  \item $\Re\rho\neq1/2$, in which case $\rho$ and $\J(\rho)$ are distinct zeros at the same height and symmetric horizontal displacement.
\end{enumerate}
Including conjugation produces a quartet unless further coincidences occur.
\end{proposition}
\begin{proof}
Immediate from the action of $\J(\sigma+it)=1-\sigma+it$ and complex conjugation.
\end{proof}

\begin{figure}[htbp]
 \centering
 \includegraphics[width=0.82\textwidth]{involution_geometry.png}
 \caption{The zeta symmetries produce an orbit. Symmetry alone permits an off-axis quartet; the Riemann Hypothesis asserts collapse onto the fixed line.}
 \label{fig:involution-geometry}
\end{figure}

This proposition isolates the missing ingredient. A proof needs a principle that forbids distinct paired points. Positivity, self-adjointness, monotonicity, or an exact cancellation law could supply such a principle, but invariance alone cannot.

\section{The affine origin of one half}

The number $1/2$ is forced by the complement map $C(\sigma)=1-\sigma$.

\begin{lemma}[Unique complement fixed point]
On $\R$, the equation $C(\sigma)=\sigma$ has the unique solution $\sigma=1/2$.
\end{lemma}

More generally, an affine involution $C_a(x)=a-x$ has centre $a/2$. The half is not numerology; it is the algebraic centre of the unit complement. The unit is supplied by the zeta functional equation, which pairs $s$ with $1-s$.

The resulting exact hierarchy is:
\begin{enumerate}
  \item the number $1$ enters through the analytic functional equation;
  \item complementarity relative to $1$ selects the centre $1/2$;
  \item the Riemann Hypothesis asks whether all non-trivial zeros lie on the associated fixed line.
\end{enumerate}

\section{Centred displacement}

Define the horizontal displacement
\begin{equation}
 x=\sigma-\frac12.
 \label{eq:x-displacement}
\end{equation}
Then complementarity acts as $x\mapsto-x$. Any scalar diagnostic intended to penalize departure from the axis should naturally be even in $x$. Binary entropy deficit, cancellation intensity near its minimum, and many positivity candidates have exactly this property.

The centred complex coordinate
\[
 u=s-\frac12=x+it
\]
transforms as
\[
 u\mapsto-\overline{u}
\]
under $\J$. In this centered vector coordinate the action is genuinely anti-linear. Multiplying by $-i$ gives the spectral coordinate $z=-iu$, under which $\J$ becomes ordinary conjugation:
\[
 z\mapsto\overline{z}.
\]
Thus a self-adjoint spectral realization would convert the centered anti-linear action into the standard reality condition for spectra.

\section{Self-duality versus zerohood}

The fixed-line property is a property of the argument $s$; zerohood is a property of the value $\xi(s)$. A canonical bridge must relate them without assuming either. The following diagram summarizes the required logic of the initial information-spinor ansatz:
\[
 \xi(s)=0
 \quad\Longrightarrow\quad
 \text{canonical state cancellation}
 \quad\Longrightarrow\quad
 \J(s)=s
 \quad\Longrightarrow\quad
 \Re s=\frac12.
\]
Only the last implication is elementary. The middle implication is the two-channel lemma developed later. The first is SOH-C001, which remains open for that ansatz.

\begin{warning}
Replacing the first arrow by a definition does not solve the problem. For example, assigning the state $\ket{-}$ whenever $\xi(s)=0$ is tautological. The state must be constructed independently and its readout zero must be proved equivalent to $\xi(s)=0$.
\end{warning}

\section{A symmetry-compatible channel swap}

Let $X$ be the swap operator on $\C^2$,
\[
 X\ket{0}=\ket{1},\qquad X\ket{1}=\ket{0}.
\]
A state map $s\mapsto\ket{\Psi(s)}$ with channel weights $(\sigma,1-\sigma)$ is naturally compatible with $\J$ if
\begin{equation}
 \ket{\Psi(\J s)}\sim X\,\overline{\ket{\Psi(s)}},
 \label{eq:state-covariance-preview}
\end{equation}
where $\sim$ denotes equality of rays. At a fixed point $s\in\Crit$, this covariance enforces equal channel norms. It does not by itself enforce the relative phase $\pi$ needed for destructive cancellation. This is the state-space version of the same logical gap: symmetry gives balance, while zerohood requires a further phase or orthogonality condition.

\status{SOH-L003 and the orbit analysis are exact. The affine map $\J(s)=1-\bar s$ is anti-holomorphic and conjugate-affine; only its centered action $u\mapsto-\bar u$ is anti-linear. The state covariance in \cref{eq:state-covariance-preview} is a design requirement for the initial candidate map, not an established property of zeta zeros.}
